% 调用上级目录公共模板
\documentclass{../template/softeng}

\begin{document}

% 封面：文档名、项目名、版本、编写人、审核人、审批人、日期
\doccover{概要设计说明书}
{在线OJ在线判题系统}
{V2026.07}

% 目录
\tableofcontents
\newpage

% 修订记录
\begin{revrecord}
V1.0 & 2026-07-11 & 张三 & 李四 & 初稿完成，完成系统总体架构、模块、接口、数据结构概要设计梳理 \\
\hline
V1.1 & 2026-07-12 & 张三 & 王五 & 补充判题流程、异常出错处理、安全与维护相关设计内容 \\
\hline
\end{revrecord}

\section{引言}
\subsection{编写目的}
本文档为在线OJ在线判题系统概要设计说明书，基于前期软件需求规格说明书完成系统顶层架构设计。文档用于明确系统整体模块划分、业务处理流程、内外接口规范、数据库数据结构、运行控制方案、异常处理、安全及维护设计，为后续详细设计、编码开发、单元测试提供设计依据。
本文档阅读对象包含项目负责人、系统架构设计师、前后端开发工程师、测试工程师。

\subsection{项目背景}
本系统为高校大二软件工程实训开发项目，开发主体为校内实训开发小组，服务于高校程序设计课程教师与在校学生。系统独立部署，无第三方业务系统对接，依靠Web服务、判题沙箱服务、关系型数据库实现完整在线编程、自动评测、教学管理业务闭环，用于支撑线上编程练习、作业自动批改、学情数据统计。

\subsection{定义}
\begin{itemize}
    \item OJ：Online Judge，在线程序自动判题系统；
    \item AC：Accepted，代码通过全部测试用例；
    \item 判题沙箱：隔离执行用户代码的安全运行环境，限制资源访问，防范恶意代码破坏服务；
    \item 测试用例：包含标准输入、预期输出、边界条件的程序评测数据集；
    \item 判题任务队列：缓存待执行判题任务的异步调度队列；
    \item RESTful API：前后端交互标准接口规范；
\end{itemize}

\subsection{参考资料}
\begin{enumerate}
    \item GB/T 9386-2008 计算机软件概要设计说明书规范；
    \item 《在线OJ在线判题系统软件需求规格说明书》V2026.07；
    \item 2026软件工程实训项目任务书；
    \item Web前后端分离开发设计规范；
    \item Linux沙箱安全隔离开发技术文档；
    \item MySQL数据库设计开发规范。
\end{enumerate}

\section{任务概述}
\subsection{目标}
搭建一套前后端分离架构的在线编程自动判题平台，划分学生、教师两类角色；拆解系统为用户模块、题库模块、在线代码编辑与判题模块、提交记录统计模块、讨论区模块、系统监控管理模块；实现题库分级管理、在线代码提交、沙箱自动编译评测、刷题数据排行、师生交流讨论、教师后台运维监控全套能力，满足高校编程课程线上练习、自动批改、学情统计教学需求。

\subsection{运行环境}
\begin{itemize}
    \item 服务端操作系统：Linux CentOS 7 / Windows Server 2019；
    \item Web服务容器：Nginx、Tomcat；
    \item 数据库：MySQL 8.0；
    \item 判题依赖：GCC、Python、Java等多语言编译器、进程隔离沙箱程序；
    \item 客户端环境：Chrome、Edge、Firefox主流现代浏览器；
    \item 中间件：Redis（缓存、任务队列）。
\end{itemize}

\subsection{需求概述}
系统分为学生端与教师管理端两大业务域：
1. 学生端：账号注册登录、分级题库浏览、在线代码编写提交、自动判题结果反馈、个人提交记录与刷题统计、全站排行榜、题目讨论区发帖交流；
2. 教师端：题库增删改查、批量导入题目、测试用例配置、全量用户与提交行为监管、异常抄袭提交标记、平台运行状态可视化监控、公告发布、讨论区内容审核管理、教学统计报表导出。

\subsection{条件与限制}
\begin{enumerate}
    \item 判题模块必须采用沙箱隔离机制，禁止用户代码直接读写服务器文件、访问网络；
    \item 单份代码运行设置时间、内存双重上限，防止服务器资源耗尽；
    \item 系统仅开放学生、教师两类角色，无游客完整答题权限；
    \item 题库批量导入仅支持标准结构化文本文件，不支持自定义复杂格式；
    \item 服务器硬件资源有限，并发判题任务存在最大并发阈值；
    \item 所有接口需携带身份鉴权凭证，严格区分角色数据访问权限。
\end{enumerate}

\section{总体设计}
\subsection{处理流程}
1. 用户注册/登录鉴权，校验身份后跳转对应功能页面；
2. 学生浏览题库，选定题目进入在线代码编辑器，编写代码提交至后端；
3. 后端接收代码、题目ID，封装判题任务送入Redis任务队列；
4. 判题服务消费队列任务，调度沙箱隔离运行代码，加载对应测试用例逐组比对输出；
5. 收集判题状态、运行耗时、内存占用结果，写入数据库提交记录；
6. 前端轮询获取判题结果，展示评测状态、错误日志；
7. 系统同步更新个人解题统计、全站排行榜数据；
8. 教师登录后台，完成题库维护、用户监管、平台监控、讨论区管理操作；
9. 所有操作行为记录系统日志，异常数据触发预警提示。

\subsection{总体结构和模块外部设计}
系统整体采用分层架构，自上而下分为四层：前端展示层、应用服务层、判题服务层、数据持久层，拆分六大独立业务模块：
\begin{enumerate}
    \item 用户权限模块：账号注册登录、身份鉴权、角色权限控制、个人信息管理；
    \item 题库管理模块：题目CRUD、难度标签管理、测试用例配置、批量导入导出；
    \item 在线判题模块：在线代码编辑器、任务队列调度、沙箱编译运行、结果判定；
    \item 数据统计排行模块：个人提交记录、解题统计、全站实时排行榜、教学报表；
    \item 讨论与公告模块：题目讨论帖子、评论回复、全站公告发布、内容审核；
    \item 系统运维监控模块：在线人数、判题队列、编译器负载、服务器健康监控、异常提交标记。
\end{enumerate}

\subsection{功能分配}
\begin{itemize}
    \item 用户权限模块：承担登录校验、账号管理、接口鉴权全部功能；
    \item 题库管理模块：负责题目基础数据、测试用例、难度分类的存储与维护；
    \item 在线判题模块：承载代码接收、异步调度、沙箱执行、结果比对核心逻辑；
    \item 数据统计排行模块：读取提交记录，实时计算通过率、解题数量、排名数据；
    \item 讨论与公告模块：管理帖子、评论、公告内容及教师审核操作；
    \item 系统运维监控模块：采集服务器与业务运行指标，提供可视化监控视图与违规标记功能。
\end{itemize}

\section{接口设计}
\subsection{外部接口}
1. 用户界面接口：Web前端可视化页面，包含题库列表、代码编辑器、监控图表、表格列表、弹窗操作提示，通过浏览器与人交互；
2. 软件外部接口：MySQL数据库连接接口、Redis缓存与队列接口、操作系统编译器调用接口；
3. 硬件接口：无专用外设接口，仅调用服务器CPU、内存、磁盘基础硬件资源。

\subsection{内部接口}
各业务模块之间通过内部RESTful接口、本地进程调用完成通信：
\begin{enumerate}
    \item 前端 ↔ 应用服务层：HTTP/HTTPS REST API，携带Token鉴权；
    \item 应用服务层 ↔ 判题服务层：Redis队列消息通信；
    \item 判题服务层 ↔ 沙箱执行程序：本地进程调用标准输入输出通信；
    \item 各业务模块 ↔ 数据持久层：JDBC数据库交互接口；
\end{enumerate}

\section{数据结构设计}
\subsection{逻辑结构设计}
数据库逻辑实体包含用户实体、题目实体、测试用例实体、提交记录实体、讨论帖子实体、评论实体、系统公告实体、系统运行日志实体，实体间关联关系：
\begin{itemize}
    \item 用户一对多提交记录、一对多讨论帖子；
    \item 题目一对多测试用例、一对多提交记录、一对多讨论帖子；
    \item 帖子一对多评论；
\end{itemize}

\subsection{物理结构设计}
采用MySQL关系型数据表存储，每张数据表设置自增主键ID、创建时间、更新时间通用字段；高频访问的排行榜、在线用户数据存入Redis内存缓存；判题任务临时数据存放Redis队列；历史提交日志采用分表存储优化查询性能。

\subsection{数据结构与程序的关系}
\begin{itemize}
    \item 用户表：由用户权限模块读写；
    \item 题目、测试用例表：由题库管理模块读写，判题模块只读；
    \item 提交记录表：判题模块写入，统计排行模块读取；
    \item 帖子、评论、公告表：讨论公告模块读写；
    \item 系统日志、监控指标表：运维监控模块读写。
\end{itemize}

\section{运行设计}
\subsection{运行模块的组合}
1. 学生正常刷题流程：用户权限模块 + 题库模块 + 在线判题模块 + 统计排行模块；
2. 教师题库管理流程：用户权限模块 + 题库管理模块；
3. 教师平台监控流程：用户权限模块 + 系统运维监控模块；
4. 师生交流流程：用户权限模块 + 讨论与公告模块；
5. 后台定时统计任务：统计排行模块独立定时运行。

\subsection{运行控制}
1. 用户访问页面时，先由用户权限模块校验Token，无权限直接跳转登录页；
2. 代码提交后任务异步入队，判题服务轮询消费队列，任务串行隔离执行；
3. 监控模块定时采集运行指标，达到阈值触发前端预警提示；
4. 教师执行删除题目、屏蔽帖子等高危操作增加二次确认弹窗；
5. 系统定时任务凌晨自动统计每日刷题报表并入库。

\subsection{运行时间}
\begin{itemize}
    \item 前端页面接口响应：≤1.5s；
    \item 单任务判题执行耗时：遵循题目设定时空限制，最大不超过10s；
    \item 监控指标刷新周期：5s；
    \item 每日统计定时任务执行时长：≤30s；
    \item 题库批量导入100道题目耗时：≤10s。
\end{itemize}

\section{出错处理设计}
\subsection{出错输出信息}
1. 用户侧错误提示：账号密码错误、权限不足、代码编译失败、运行超时、内存超限、答案错误、网络连接失败；
2. 教师后台错误提示：题目导入格式错误、测试用例缺失、批量操作失败、服务器负载过高；
3. 系统后台日志错误：数据库连接失败、判题服务宕机、队列阻塞、恶意代码拦截记录。

\subsection{出错处理对策}
\begin{enumerate}
    \item 网络/数据库断开：前端缓存未提交代码，后台定时重试数据库连接，页面提示稍后访问；
    \item 判题服务宕机：未处理任务留存队列，服务恢复后自动续判，前端提示判题服务繁忙；
    \item 恶意超限代码：沙箱强制终止进程，直接返回运行超限结果，记录违规日志；
    \item 高并发队列阻塞：限制新提交入口，提示当前服务器繁忙，请稍后再提交；
    \item 数据操作异常：开启数据库事务，出错自动回滚，避免脏数据入库。
\end{enumerate}

\section{安全保密设计}
\begin{enumerate}
    \item 身份安全：采用Token登录鉴权，密码加密存储，接口访问权限按角色隔离；
    \item 代码运行安全：沙箱隔离用户代码，禁用文件读写、网络调用、系统高危指令；
    \item 接口安全：请求参数过滤，防御SQL注入、XSS跨站脚本攻击；
    \item 数据保密：学生提交代码、个人刷题记录仅本人与授课教师可见；
    \item 操作审计：所有后台修改、删除操作完整记录操作人、操作时间、操作内容日志。
\end{enumerate}

\section{维护设计}
\begin{enumerate}
    \item 日志维护模块：统一收集访问日志、判题日志、异常日志，支持按时间、用户检索；
    \item 监控运维面板：可视化展示服务器负载、判题队列、在线人数，便于运维定位故障；
    \item 模块化分层设计：各业务模块低耦合，可单独迭代、替换升级；
    \item 数据备份机制：定时自动备份MySQL题库、提交记录核心业务数据；
    \item 配置文件分离：服务器地址、判题并发数、时空限制等参数外置配置，无需修改代码即可调整。
\end{enumerate}

% 插入系统总体架构图示例
\begin{figure}[htbp]
    \centering
    % \includegraphics[width=0.7\textwidth]{oj_arch.png}
    \caption{在线OJ系统总体架构图}
\end{figure}

% 代码块示例：判题任务队列实体结构
\begin{lstlisting}[language=Python]
import dataclasses

@dataclasses.dataclass
class JudgeTask:
    submit_id: int       # 提交记录ID
    user_id: int         # 提交学生ID
    problem_id: int      # 题目ID
    code: str            # 用户源代码
    lang: str            # 编程语言
    time_limit: int      # 时间限制 ms
    memory_limit: int    # 内存限制 KB
\end{lstlisting}

\end{document}